% from ega3-6-fr.tex
\oldpage[III]{147}
S-morphismes $u:X^{(1)}\to X$, $v:Y^{(1)}\to Y$, correspondant à des
$A$-homomorphismes $\varphi:B\to B^{(1)}$, $\psi:C\to C^{(1)}$.
Considérons les complexes
$u^*(\mathscr P_\bullet)=(P_\bullet\otimes_B B^{(1)})^{\sim}$ de
$\mathscr O_{X^{(1)}}$-Modules,
$v^*(\mathscr Q_\bullet)=(Q_\bullet\otimes_C C^{(1)})^{\sim}$ de
$\mathscr O_{Y^{(1)}}$-Modules. Les $A$-homomorphismes canoniques
\[
  P_\bullet\longrightarrow P_\bullet\otimes_B B^{(1)},
  \qquad
  Q_\bullet\longrightarrow Q_\bullet\otimes_C C^{(1)}
\]
donnent par fonctorialité un $A$-homomorphisme
\[
  \mathbf{Tor}_\bullet^A(P_\bullet,Q_\bullet)
  \longrightarrow
  \mathbf{Tor}_\bullet^A
  (P_\bullet\otimes_B B^{(1)},Q_\bullet\otimes_C C^{(1)});
\]
en outre, toujours par fonctorialité, cet homomorphisme est en fait un
homomorphisme de $(B\otimes_A C)$-modules. D'où l'on conclut que l'on a défini
ainsi un $(u\times_S v)$-morphisme
\[
  \theta:\mathscr{T}\!or_\bullet^S(\mathscr P_\bullet,\mathscr Q_\bullet)
  \longrightarrow
  \mathscr{T}\!or_\bullet^S
  (u^*(\mathscr P_\bullet),v^*(\mathscr Q_\bullet))
  \tag{6.4.4.1}\label{III.6.4.4.1.fr}
\]
et par suite, un homomorphisme de
$\mathscr O_{X^{(1)}\times_S Y^{(1)}}$-Modules
\[
  \theta^\sharp:(u\times_S v)^*
  \bigl(\mathscr{T}\!or_\bullet^S(\mathscr P_\bullet,\mathscr Q_\bullet)\bigr)
  \longrightarrow
  \mathscr{T}\!or_\bullet^S
  (u^*(\mathscr P_\bullet),v^*(\mathscr Q_\bullet))
  \tag{6.4.4.2}\label{III.6.4.4.2.fr}
\]
qui est évidemment un morphisme de bi-$\partial$-foncteurs dans les catégories
de Modules quasi-cohérents limités inférieurement.

L'homomorphisme (6.4.4.2) n'est pas nécessairement bijectif; toutefois :
\end{env}

\begin{namedtheorem}[6.4.5]{Lemme}
Avec les notations de (6.4.4), supposons que $u$ et $v$ soient des immersions
ouvertes; alors l'homomorphisme (6.4.4.2) est bijectif.
\end{namedtheorem}

\begin{proof}
Identifions $X^{(1)}$ (resp. $Y^{(1)}$) à un ouvert de $X$ (resp. $Y$);
$X^{(1)}$ (resp. $Y^{(1)}$) est alors réunion d'ouverts de la forme
$\mathrm D(f)$ (resp. $\mathrm D(g)$), où $f\in B$ (resp. $g\in C$), et les
préschémas induits $\mathrm D(\varphi(f))$ et $\mathrm D(f)$ (resp.
$\mathrm D(\psi(g))$ et $\mathrm D(g)$) sont isomorphes. Il suffira de prouver
le lemme lorsque $X^{(1)}$ (resp. $Y^{(1)}$) est de la forme $\mathrm D(f)$
(resp. $\mathrm D(g)$); en effet, si ce point est établi, et si l'on revient au
cas général, il suffira de prouver que la restriction de $\theta^\sharp$ à
chaque ouvert $\mathrm D(f)\times_S\mathrm D(g)$ est un isomorphisme; or, si
$u_1:\mathrm D(f)\to X^{(1)}$, $v_1:\mathrm D(g)\to Y^{(1)}$ sont les
injections canoniques, la restriction précédente n'est autre que
$(u_1\times_S v_1)^*(\theta^\sharp)$; mais il est immédiat, en vertu des
définitions (6.4.4) et de ($0_{\mathrm I}$, 4.4.8), qu'en la composant avec
l'homomorphisme canonique
\[
  (u_1\times_S v_1)^*
  \bigl(\mathscr{T}\!or_\bullet^S
    (u^*(\mathscr P_\bullet),v^*(\mathscr Q_\bullet))\bigr)
  \longrightarrow
  \mathscr{T}\!or_\bullet^S
  ((u')^*(\mathscr P_\bullet),(v')^*(\mathscr Q_\bullet))
  \tag{6.4.5.1}\label{III.6.4.5.1.fr}
\]
où $u'=u\circ u_1$ et $v'=v\circ v_1$, on obtient l'homomorphisme canonique
\[
  (u'\times_S v')^*
  \bigl(\mathscr{T}\!or_\bullet^S(\mathscr P_\bullet,\mathscr Q_\bullet)\bigr)
  \longrightarrow
  \mathscr{T}\!or_\bullet^S
  ((u')^*(\mathscr P_\bullet),(v')^*(\mathscr Q_\bullet))
  \tag{6.4.5.2}\label{III.6.4.5.2.fr}
\]
et si l'on sait que (6.4.5.1) et (6.4.5.2) sont des isomorphismes, il en
résultera qu'il en est de même de $(u_1\times_S v_1)^*(\theta^\sharp)$.

Supposons donc que $X^{(1)}=\mathrm D(f)$ et $Y^{(1)}=\mathrm D(g)$, de sorte
que $B^{(1)}=B_f$ et $C^{(1)}=C_g$; $u^*(\mathscr P_\bullet)$ (resp.
$v^*(\mathscr Q_\bullet)$) s'identifie alors à
$((P_\bullet)_f)^{\sim}$ (resp. $((Q_\bullet)_g)^{\sim}$); d'autre part,
$X^{(1)}\times_S Y^{(1)}$ s'identifie au sous-schéma ouvert
$\mathrm D(f\otimes g)$ de
$X\times_S Y=\operatorname{Spec}(B\otimes_A C)$ (II, 4.3.2.4); il s'agit de
prouver que l'homomorphisme
\[
  \bigl(\mathbf{Tor}_\bullet^A(P_\bullet,Q_\bullet)\bigr)_{f\otimes g}
  \longrightarrow
  \mathbf{Tor}_\bullet^A((P_\bullet)_f,(Q_\bullet)_g)
  \tag{6.4.5.3}\label{III.6.4.5.3.fr}
\]

